\chapter{DRIVE KINEMATICS AND WHEEL ODOMETRY}
\label{ch:kinematics}

\section{NOTATION}

Table~\ref{tab:geom} lists the geometric parameters, with the values in
\file{ros/bbai\_base/config/base.yaml}.

\begin{table}[htbp]
\centering
\caption{Geometry and drive parameters}
\label{tab:geom}
\small
\begin{tabular}{clll}
\toprule
Symbol & Meaning & Value & \texttt{base.yaml} key \\
\midrule
$r$ & wheel radius & \SI{0.0396}{m} & \texttt{wheel\_radius} \\
$b$ & track (left--right wheel centres) & \SI{0.120}{m} & \texttt{track\_width} \\
$L$ & wheelbase (front--back axles) & \SI{0.101}{m} & --- \\
$N$ & encoder counts per wheel revolution & \num{15} & \texttt{counts\_per\_rev} \\
$\Delta t$ & odometry period & $1/30$~s & \texttt{rate} \\
$v_{\max}$ & wheel speed at duty 1 (lifted) & \SI{1.2}{m/s} & \texttt{max\_wheel\_speed} \\
$d_{\max}$ & duty limit & 0.60 & \texttt{max\_duty} \\
$d_{\min}$ & duty floor for any nonzero command & 0.12 & \texttt{min\_duty} \\
$T_{\mathrm{cmd}}$ & command timeout & \SI{0.5}{s} & \texttt{cmd\_timeout} \\
\bottomrule
\end{tabular}
\end{table}

The pose in the \frm{odom} frame is $\bm{p}=[x\ y\ \psi]\T$. Subscript $k$
denotes the $k$-th cycle of the \SI{30}{Hz} loop, and $\Delta(\cdot)_k$ the
change between cycles $k-1$ and $k$.

\section{FROM ENCODER COUNTS TO WHEEL TRAVEL}

Let $c_j$ be the raw count of encoder channel $j\in\{1,2,3,4\}$ and
$\sigma_j\in\{+1,-1\}$ its sign from Table~\ref{tab:channels}. One count is an
arc of
\begin{equation}\label{eq:mpc}
\delta = \frac{2\pi r}{N} = \frac{2\pi(0.0396)}{15} = \SI{0.01659}{m}.
\end{equation}
Each side averages its two wheels,
\begin{equation}
s_L = \frac{\delta}{2}\bigl(\sigma_1c_1+\sigma_3c_3\bigr),\qquad
s_R = \frac{\delta}{2}\bigl(\sigma_2c_2+\sigma_4c_4\bigr),
\end{equation}
and the travel of the body centre in one cycle is
\begin{equation}\label{eq:dcentre}
d_k = \tfrac12\bigl(\Delta s_{L,k}+\Delta s_{R,k}\bigr).
\end{equation}
Averaging both wheels per side halves the effect of a single wheel slipping
or being spun by hand: a spurious count on one wheel moves its side's average
by half a count.

\begin{remark}[Quantization]\label{rem:quant}
With $N=15$ the encoders are coarse. One count is \SI{16.6}{mm} of wheel
travel, so $d_k$ is quantized in steps of $\delta/4=\SI{4.1}{mm}$. The
velocity published on \topic{/odom}, $v_k=d_k/\Delta t$, is therefore
quantized in steps of $\delta/(4\Delta t)=\SI{0.124}{m/s}$ at \SI{30}{Hz}: at
a steady \SI{0.3}{m/s} it alternates between about 2 and 3 levels from one
message to the next. This is the same trade the thesis identifies for its
motor encoders (the sample-rate remark in thesis \S7.5): raising the loop rate makes the
differenced velocity noisier for a fixed encoder resolution. Position, which
is a sum of counts, is not affected; only the per-message velocity is.
Anything that consumes the twist (the EKF in particular) should either use a
longer differencing window or carry a velocity variance of at least
$\delta^2/(192\,\Delta t^2)\approx\SI{1.3e-3}{m^2/s^2}$, the variance of a
uniform quantization error of width $\delta/(4\Delta t)$.
\end{remark}

\section{SKID-STEER KINEMATICS}
\label{sec:skid}

For a differential-drive robot with no wheel slip the body speed and yaw rate
are
\begin{equation}\label{eq:diffdrive}
v = \tfrac12(v_R+v_L),\qquad \omega = \frac{v_R-v_L}{b}.
\end{equation}
A four-wheel skid-steer robot cannot turn without its wheels sliding
sideways, because the front and back wheels on each side are fixed parallel
at a distance $L$ apart. The instantaneous centres of rotation of the left and
right treads move outward, which is captured by an effective track
$\chi b$ with $\chi\ge1$ \cite{mandow2007skid,martinez2005icr}:
\begin{equation}\label{eq:skidkin}
\omega = \frac{v_R-v_L}{\chi\,b}.
\end{equation}
$\chi$ depends on the ratio $L/b$ (here $0.84$, which is high), the floor
surface, the load and the turn rate itself; values of 1.5 to 2 are typical. An
encoder-only heading using \eqref{eq:diffdrive} would overestimate every turn
by the factor $\chi$, and that factor would change between carpet and tile.
This is why the base node takes its heading from the gyroscope
(\texttt{use\_gyro\_heading: true}) and uses the encoders only for distance.

\section{ODOMETRY INTEGRATION}
\label{sec:odom}

The heading increment is
\begin{equation}\label{eq:dpsi}
\Delta\psi_k = \begin{cases}
\bigl(g_{z,k}-\hat b_g\bigr)\Delta t_k & \text{gyro heading (default)}\\[2pt]
\dfrac{\Delta s_{R,k}-\Delta s_{L,k}}{b} & \text{encoder heading}
\end{cases}
\end{equation}
where $g_z$ is the $z$ gyroscope rate and $\hat b_g$ the bias estimate
(\S\ref{sec:gyro}). While the robot is parked, $\Delta\psi_k=0$
(\S\ref{sec:zupt}). The pose is advanced with the midpoint rule
(second-order Runge--Kutta),
\begin{equation}\label{eq:odomint}
\begin{aligned}
x_{k} &= x_{k-1} + d_k\cos\!\left(\psi_{k-1}+\tfrac12\Delta\psi_k\right),\\
y_{k} &= y_{k-1} + d_k\sin\!\left(\psi_{k-1}+\tfrac12\Delta\psi_k\right),\\
\psi_{k} &= \wrap\!\left(\psi_{k-1}+\Delta\psi_k\right),
\end{aligned}
\end{equation}
with $\wrap(\alpha)=\atantwo(\sin\alpha,\cos\alpha)\in(-\pi,\pi]$.
For motion on a circular arc the exact update is
\begin{equation}
x_k = x_{k-1}+\frac{d_k}{\Delta\psi_k}\Bigl[\sin(\psi_{k-1}+\Delta\psi_k)-\sin\psi_{k-1}\Bigr],
\end{equation}
and similarly for $y$. Expanding both in $\Delta\psi_k$ shows that the
midpoint rule differs from the exact arc only by a factor
$1-\Delta\psi_k^2/24$ on the step length. At the fastest turn the robot makes,
\SI{10}{rad/s}, $\Delta\psi_k=\SI{0.33}{rad}$ and the error is 0.5\% of
$d_k$, and in a pure spin $d_k\approx0$ anyway, so \eqref{eq:odomint} is
adequate.

The message also carries the body velocity $v_k=d_k/\Delta t_k$ and
$\omega_k=\Delta\psi_k/\Delta t_k$, and the same pose is broadcast as the
transform \frm{odom}$\to$\frm{base\_link} with quaternion
$\bm{q}=[0,\,0,\,\sin\frac{\psi}{2},\,\cos\frac{\psi}{2}]$.

\subsection{Odometry uncertainty}

Linearising \eqref{eq:odomint} about the current estimate gives the usual
dead-reckoning covariance propagation \cite{thrun2005probabilistic},
\begin{equation}\label{eq:odomcov}
P_k = F_k P_{k-1}F_k\T + G_k Q_k G_k\T,\qquad
F_k=\begin{bmatrix}1&0&-d_k\sin\bar\psi\\0&1&d_k\cos\bar\psi\\0&0&1\end{bmatrix},\quad
G_k=\begin{bmatrix}\cos\bar\psi&0\\\sin\bar\psi&0\\0&1\end{bmatrix},
\end{equation}
where $\bar\psi=\psi_{k-1}+\frac12\Delta\psi_k$ and
$Q_k=\operatorname{diag}(\sigma_d^2,\sigma_{\psi}^2)$ is the variance of one
step's distance and heading increments. The position uncertainty grows
without bound, and the heading term couples into it through the off-diagonal
entries of $F_k$: a heading error $\delta\psi$ produces a lateral error
$d\,\delta\psi$ after travelling $d$.

\begin{remark}[The published covariance is constant]
\texttt{base\_node.py} publishes a fixed diagonal covariance,
$(0.01, 0.01, 10^6, 10^6, 10^6, 0.05)$ for $(x,y,z,\phi,\theta,\psi)$ on both
the pose and the twist. The $10^6$ entries correctly mark the unobserved
states. The pose entries, however, claim a \SI{0.1}{m} standard deviation
forever, while \eqref{eq:odomcov} grows. A filter that fuses the
\topic{/odom} \emph{pose} as an absolute measurement will therefore trust it
far too much after a few metres. The safe use is to fuse only the twist, or the
pose in robot\_localization's \texttt{differential} mode. The GPS EKF does the
former: it takes only $v_x$ from \topic{/odom} (\S\ref{sec:ekfcfg}).
\end{remark}

\section{FROM VELOCITY COMMAND TO MOTOR DUTY}
\label{sec:duty}

A \topic{/cmd\_vel} message $(v,\omega)$ is turned into a desired speed for each
side using the skid-steer model \eqref{eq:skidkin} with a commanded
effective track $\chi_c$,
\begin{equation}\label{eq:sidespeed}
v_L = v-\tfrac12\omega \chi_c,\qquad v_R = v+\tfrac12\omega \chi_c,
\end{equation}
and each side speed into a PWM duty for both of that side's motors,
\begin{equation}\label{eq:duty}
u=\frac{v_{\mathrm{side}}}{v_{\max}},\qquad
D = \sat_{d_{\max}}\!\Bigl(\sgn(u)\bigl[d_{\min}+(1-d_{\min})\abs{u}\bigr]\Bigr),
\quad (D=0 \text{ if } \abs{u}<10^{-3}),
\end{equation}
multiplied by the motor's sign.

\paragraph{Effective track for commands.} Originally $\chi_c$ was the
geometric track $b=\SI{0.12}{m}$. On the floor a command of
\SI{5.0}{rad/s} then produced only \SI{1.31}{rad/s} as measured by the gyro,
because the tyres scrub sideways and the robot turns as if its track were
wider (\S\ref{sec:skid}). The ratio gives the effective track directly,
\begin{equation}\label{eq:chicmd}
\chi_c = b\,\frac{\omega_{\mathrm{cmd}}}{\omega_{\mathrm{meas}}}
       = 0.12\times\frac{5.0}{1.31}\approx\SI{0.46}{m},
\end{equation}
now the \texttt{cmd\_turn\_track} parameter in \file{base.yaml}, so that
\topic{/cmd\_vel}'s $\omega$ approximately equals the real yaw rate on that
floor. Odometry is unaffected: it still takes heading from the gyro.
\influx{\texttt{cmd\_turn\_track} is on branch \texttt{slam-estimator-updates}
(commit 2d2a56d), not yet on \texttt{main}. It was measured on one floor;
carpet and hard floor will differ.}

The offset $d_{\min}$ is a deadzone
compensator: below about 8\% duty the wheels do not turn at all (lifted
test), so any nonzero command starts at 12\%. This is the inverse of the
deadzone nonlinearity the thesis models for its own motors (thesis
\S7.6.2). The control is open loop: duty sets voltage, and the speed that
results depends on load and battery voltage. The encoders are used for
odometry but not for speed feedback.

Table~\ref{tab:dutyex} evaluates \eqref{eq:sidespeed}--\eqref{eq:duty} for
the commands teleop sends, with the old and new $\chi_c$. The duty limit of
0.6 is reached at $\abs{u}=0.545$, so the fastest either side can be
commanded is about \SI{0.65}{m/s}, and with $\chi_c=\SI{0.46}{m}$ the
fastest commanded spin is $2\times0.65/0.46\approx\SI{2.8}{rad/s}$.

\begin{table}[htbp]
\centering
\caption{Commanded duty $(D_L, D_R)$ for typical commands}
\label{tab:dutyex}
\small
\begin{tabular}{lcc}
\toprule
Command & $\chi_c=b=\SI{0.12}{m}$ (old) & $\chi_c=\SI{0.46}{m}$ (new) \\
\midrule
Forward \SI{0.3}{m/s} & $0.34,\ 0.34$ & $0.34,\ 0.34$ \\
Spin \SI{1}{rad/s} & $-0.16,\ 0.16$ & $-0.29,\ 0.29$ \\
Spin \SI{1.5}{rad/s} (old full stick) & $-0.19,\ 0.19$ & $-0.37,\ 0.37$ \\
Spin \SI{6}{rad/s} (new full stick) & $-0.38,\ 0.38$ & $-0.60,\ 0.60$ (limit) \\
Spin \SI{10}{rad/s} (new turbo) & $-0.56,\ 0.56$ & $-0.60,\ 0.60$ (limit) \\
Forward \SI{0.3}{m/s} + \SI{1.5}{rad/s} & $0.27,\ 0.41$ & $-0.15,\ 0.59$ \\
\bottomrule
\end{tabular}
\end{table}

\section{WHY SKID-STEER TURNING NEEDS SO MUCH POWER}
\label{sec:skidturn}

Jason observed on 2 October 2026 that normal-speed turning would not start the
robot spinning from rest and barely turned it while moving. A simple
friction model explains this \cite{wong2008theory}. Let the robot weigh $W$,
shared equally by the four wheels, which sit at $\pm L/2$ from the centre
along the body and $\pm b/2$ across it. Spinning in place, each wheel must
slide sideways, and the floor resists with a lateral friction force
$\mu_\ell W/4$ acting at a lever arm of $L/2$. The total resisting moment is
\begin{equation}
M_r = 4\cdot\mu_\ell\frac{W}{4}\cdot\frac{L}{2}=\frac{\mu_\ell W L}{2}.
\end{equation}
The motors supply it with opposing tractive forces $\pm F$ per side at lever
arm $b/2$, a moment of $F\,b$. Turning starts only when
\begin{equation}\label{eq:spinforce}
F \ \ge\ \frac{\mu_\ell L}{2b}\,W \approx 0.42\,\mu_\ell W .
\end{equation}
Driving straight only has to overcome rolling resistance $f_rW$ with $f_r$ of
order 0.02--0.05, whereas lateral sliding friction $\mu_\ell$ is of order 0.3--0.7
on household floors. A spin therefore needs roughly 3 to 15 times the
tractive force of driving straight. With the old linear mapping, full stick
asked for \SI{1.5}{rad/s}, which with the geometric track is a duty of only
0.19 (Table~\ref{tab:dutyex}), barely above the 0.12 floor; that is enough
to keep a lifted wheel turning but not to make the tyres scrub. This
motivated the new teleop turn curve (\SI{6}{rad/s} at full stick) and then
the effective command track $\chi_c=\SI{0.46}{m}$, which roughly doubles the
duty of every turn command. A closed-loop yaw-rate or wheel-speed controller
(Chapter~\ref{ch:motor}) would make turn response independent of the floor
surface.

\begin{remark}
\eqref{eq:spinforce} shows the turn force grows with $L/b$. Moving the
wheels farther apart, or using wheels with less lateral grip, would reduce it.
\end{remark}
